\section{Three regimes}\label{sec:regimes}

What reaches a pixel depends on how the distance $r$ to the source compares
with the wavelength $\lambda = 2\pi c/\omega$ of the radiation it would emit
at angular frequency $\omega$. Table~\ref{tab:regimes} summarises the three
cases.

\begin{table}
\begin{tabular}{lll}
Regime & What the pixel receives & What limits the image \\ \hline
Near field & the changing Newtonian pull of nearby mass & blur grows with distance \\
Waves from bulk motion & radiation from accelerating bodies & each body is an unresolved point \\
Thermal glow & radiation from molecular motion & almost no gravitons arrive \\
\end{tabular}
\caption{The three regimes of a gravitational camera.}\label{tab:regimes}
\end{table}

\paragraph{Near field.} For $r \ll \lambda/2\pi$ the field is that of
Eq.~\eqref{eq:newton}, following the source almost instantaneously. At
$1\,\mathrm{Hz}$ this zone extends to about $48{,}000\,\mathrm{km}$, so every
everyday scene lies deep inside it. The near field is strong and passes through
all matter, but it obeys Laplace's equation, and a harmonic field loses detail
with distance: a pattern of spatial wavenumber $k$ in a source at depth $z$ is
weakened by $e^{-kz}$ at the camera. With a signal-to-noise ratio $\SNR$, detail
finer than
\begin{equation}\label{eq:resolution}
\Delta x \approx \frac{\pi z}{\ln \SNR}
\end{equation}
is lost. Resolution improves only with the logarithm of sensitivity: a
signal-to-noise ratio of $10^{30}$ still leaves a blur of about three degrees.

\paragraph{Waves from bulk motion.} A body of size $L$ whose parts move at
speed $v$ radiates at frequencies near $v/L$, so that
\begin{equation}\label{eq:unresolved}
\frac{\lambda}{L} \sim \frac{c}{v} .
\end{equation}
The wavelength exceeds the body by the factor by which light outruns it. No
aperture can then resolve the body in its own radiation: each moving object is
a point, distinguished only by the strength and polarization of its waves.

\paragraph{Thermal glow.} Sharp images need wavelengths shorter than the
scene, and the only motion in ordinary matter fast enough is thermal. A
gravitational wave crossing a medium of shear viscosity $\eta$ loses energy at
the rate found by Hawking \cite{Hawking1966}, which corresponds to an
absorption coefficient
\begin{equation}\label{eq:alpha}
\alpha = \frac{16\pi G\,\eta}{c^{3}}
\end{equation}
per unit length. A medium that absorbs must also emit. With the Planck
function $B_\omega(T)$, Kirchhoff's law gives the emissivity
\begin{equation}\label{eq:kirchhoff}
j_\omega = \alpha\, B_\omega(T),
\end{equation}
which for long wavelengths is the known result that a thermal plasma radiates
gravitational waves in proportion to its shear viscosity
\cite{GhiglieriLaine2015}. For water, $\alpha \approx 10^{-37}\,\mathrm{m^{-1}}$:
matter is transparent, and its own glow is correspondingly faint. Thermal
radiation from the vibrations of a solid has been calculated before
\cite{SchaeferDehnen1981,DehnenRomeroSchaefer1984}; for an elastic sphere
without internal damping its strength was found to grow with the sphere's
surface, not its volume. How that channel compares with the volume emission of
Eq.~\eqref{eq:kirchhoff} is an open question for this work, since it decides
whether a warm object appears as a glowing volume or a glowing outline.
